\begin{textsample}{NEG Dim 1 – 2013\_c5 – Score 8.00 – 2013\_c5/enade\_2013\_med\_9c8209e0a3.txt}  \label{ex:f1_neg_356}
CAPÍTULO 5 DISTRIBUIÇÃO DOS CONCEITOS 5.1 PANORAMA NACIONAL DA DISTRIBUIÇÃO DOS CONCEITOS A Tabela 5.1 apresenta a quantidade e distribuição de cursos de Medicina participantes do ENADE/2013, por faixa de conceito e Grande Região.
A diferença entre os cursos tabulados neste capítulo e no capítulo 2 são os cursos sem conceito, em princípio, aqueles sem alunos concluintes que participassem da prova ou que tivessem tão somente um aluno concluinte13.
Observando -se os dados da Tabela 5.1, nota -se que, dos 166 cursos participantes, 59 ( 35,5 \% ) classificaram -se com conceito 3, o valor modal.
Este foi também o conceito modal nas regiões Norte ( nesta região, além do conceito modal já citado, também foi classificado como conceito modal, o conceito 2 ), Sudeste e Sul ( entre 35,7 \% na região Norte e 41,4 \% na região Sul ).
O conceito 4 foi o segundo mais frequente em nível nacional ( 28,9 \%, correspondendo a 48 cursos ) e o conceito 2, o terceiro ( 16,9 \%, correspondendo a 28 cursos ).
Houve, ainda, 20 cursos ( 12,0 \% ) que receberam conceito 5 e outros 11 cursos que receberam conceito 1.
Dos 166 cursos de Medicina, nenhum ficou sem conceito ( SC ).
A região Norte participou com 14 cursos ou 8,4 \% do total nacional.
Nesta região, os conceitos 2 e 3 ( valores modais ) foram atribuídos a cinco cursos, \textbf{cada}.
Apenas um curso ficou com conceito 4 ( 7,1 \% ), enquanto, três cursos foram avaliados com conceito 1 ( 21,4 \% ).
Nenhum curso ficou sem conceito ( SC ) ou ficou alocado ao conceito 5.
A região Nordeste participou com 34 cursos ou 20,5 \% do total nacional.
Nesta 13 Estes cursos com somente um concluinte não foram considerados no capítulo 2. a região Nordeste participou com 34 cursos ou 20,5 \% do total nacional.
Nesta região, 15 cursos, 44,1 \% em termos regionais, obtiveram conceito 4, o conceito modal para a região.
Os conceitos 2 e 3 foram atribuídos a, respectivamente, sete e nove cursos ( 20,6 \% e 26,5 \% ).
Já os conceitos 1 e 5 foram atribuídos a, respectivamente, um e dois cursos ( 2,9 \% e 5,9 \% ).
Nessa região, nenhum dos cursos ficou sem conceito ( SC ).
Como já comentado, dos 77 cursos participantes da região Sudeste, 31 ( 40,3 \% ) obtiveram conceito 3, o conceito modal.
O conceito 1 foi atribuído a seis cursos ( 7,8 \% ) e o conceito 2, a 14 ( 18,2 \% ).
O conceito 4 foi atribuído a 17 cursos ( 22,1 \% ) e nove outros ( 11,7 \% ) receberam o conceito 5.
Nenhum curso ficou sem conceito ( SC ).
A região Sul também contou com cursos distribuídos em todas as faixas de conceitos, à exceção do conceito 2, que não foi alocado a nenhum curso.
A predominância do conceito 3 foi de 41,4 \%, como já comentado anteriormente, correspondentes a 12 dos 29 cursos participantes na região Sul.
O conceito 4 foi atribuído a 11 cursos ( 37,9 \% ), enquanto o conceito 1 foi recebido por um único curso ( 3,4 \% ) e o conceito 5, por cinco deles ( 17,2 \% ).
Nenhum dos cursos da região Sul ficou sem conceito ( SC ).
A região Centro-Oeste obteve 12 cursos participantes ou 7,2 \% do total nacional.
Nesta região nenhum curso ficou sem conceito ( SC ) ou recebeu o conceito 1.
Quatro cursos receberam igualmente os conceitos 3 e 5 ( dois cursos, \textbf{cada} ), enquanto oito cursos ficaram alocados igualmente, também, nos conceitos 4 e 5 ( quatro cursos, \textbf{cada} ).
O Gráfico 5.1 apresenta a distribuição do Conceito ENADE segundo Grande Região : poligonais mais à esquerda representam Grandes Regiões com pior distribuição, e poligonais mais à direita, Grandes Regiões com cursos com melhores conceitos.
A região Norte ( linha azul ) apresentou os cursos com valores mais baixos, ao passo que as regiões Sul ( linha roxa ) e Centro-Oeste ( linha preta ) apresentam as distribuições com valores maiores.
Cursos em IES das regiões Nordeste ( linha verde ) e Sudeste ( linha vermelha ) apresentam uma situação intermediária. 5.2 CONCEITOS POR CATEGORIA ADMINISTRATIVA E POR GRANDE REGIÃO A Tabela 5.2 apresenta a distribuição dos cursos participantes do ENADE/2013 de Medicina, por Categoria Administrativa, de acordo com os conceitos por eles alcançados, segundo as Grandes Regiões brasileiras.
Dos 166 cursos participantes, 67 ( 40,4 \% ) eram ministrados em Instituições Públicas e 99 ( 59,6 \% ), em Privadas.
De acordo com as informações da Tabela 5.2, em termos nacionais, das 20 IES que receberam o conceito 5, 16 eram Públicas e quatro Privadas.
Dos 67 cursos participantes de IES Públicas, o conceito 4 foi o valor modal, atribuído a 28 cursos ( 41,8 \% ).
Entre os demais cursos de Instituições Públicas participantes, somente um obteve conceito 1 ( 1,5 \% da categoria ), dois receberam conceito 2 ( 3,0 \% ) e 20 foram alocados ao conceito 3 ( 29,9 \% ).
Como já comentado, os 16 restantes receberam o conceito 5 ( 23,9 \% ).
Nesta categoria, nenhum dos cursos ficou sem conceito ( SC ).
Na rede privada, o conceito modal foi 3, com 39 cursos dos 99 da categoria.
Entre os demais cursos participantes, dez receberam conceito 1 e 26, conceito 2.
O conceito 4 foi atribuído a 20 cursos e o conceito 5, como já comentado, a quatro cursos.
Nesta Categoria Administrativa, nenhum dos cursos ficou sem conceito ( SC ).
Na análise por região, observa -se que, na região Norte, as Instituições Privadas participaram com sete cursos ( 50,0 \% do total regional ), apenas os conceitos 1 e 2 obtiveram cursos associados nessa categoria, dos quais quatro receberam conceito 2, o conceito modal, enquanto três cursos obtiveram o conceito 1.
Nenhum ficou sem conceito ( SC ) nessa categoria.
As Instituições Públicas participaram, também, com sete cursos ( outros 50,0 \% do total regional ), dos quais nenhum ficou sem conceito ( SC ) ou recebeu os conceitos 1 ou 5.
O conceito modal para esta Categoria Administrativa, na região Norte, foi 3, com cinco cursos, correspondendo a 71,4 \%.
Os dois cursos restantes ficaram alocados igualmente nos conceitos 2 e 4 ( um curso, \textbf{cada} ) nessa categoria e região.
Na região Nordeste, a rede privada concentrou 13 dos 34 cursos participantes, equivalentes a 38,2 \% do total da região, a menor proporção entre todas as regiões.
Nesta Categoria Administrativa na região, predominaram os cursos com conceito 2, seis cursos correspondendo a 46,2 \% dos cursos oferecidos por IES Privadas no Nordeste.
Apenas um curso recebeu o conceito 1 ( 7,7 \% ), enquanto seis receberam igualmente os conceitos 3 e 4 ( três cursos, \textbf{cada} ).
Nenhum curso ficou sem conceito ( SC ) ou recebeu o conceito 5 nessa categoria e região.
As Instituições Públicas dessa região participaram com 21 cursos ( 61,8 \% ), dos quais 12 obtiveram conceito 4, o valor modal.
Os demais foram avaliados com os conceitos 2 ( um curso ), 3 ( seis cursos ) e 5 ( dois cursos ).
Nenhum curso ficou sem conceito ( SC ) ou recebeu o conceito 1 nesta combinação de categoria e região.
Na região Sudeste, a proporção de cursos da rede privada, 74,0 \%, foi mais elevada do que nas demais regiões brasileiras, correspondendo a 57 dos 77 cursos participantes.
Nesta categoria, na região Sudeste, o conceito modal foi 3 ( 26 cursos ).
Nesta combinação de Categoria Administrativa e Grande Região, nenhum dos cursos ficou sem conceito ( SC ).
Os demais foram avaliados com conceito 1 ( seis cursos ), conceito 2 ( 14 cursos ), conceito 4 ( nove cursos ) e conceito 5 ( dois cursos ).
Entre os 20 cursos em Instituições Públicas na região Sudeste, a categoria modal foi a que recebeu conceito 4 ( oito cursos ).
Os demais foram avaliados com os conceitos 3 ( cinco cursos ) e 5 ( sete cursos ).
Na região Sudeste, nesta categoria, nenhum curso ficou sem conceito ( SC ) ou recebeu os conceitos 1 ou 2.
As Instituições Privadas concentraram 16 dos 29 cursos participantes da região Sul, 55,2 \% do total regional.
Desses, oito ficaram com conceito 3, o conceito modal.
Nesta combinação de Categoria Administrativa e Grande Região, nenhum curso ficou sem receber conceito ( SC ) ou recebeu os conceitos 1 ou 2.
Os demais foram avaliados com conceito 4 ( seis cursos ) e conceito 5 ( dois cursos ).
As Instituições Públicas na região Sul participaram com 13 cursos ( 44,8 \% ), dos quais nenhum ficou sem conceito ( SC ) ou recebeu o conceito 2.
O conceito modal foi 4, com cinco cursos.
Além disso, os demais foram avaliados com conceito 1 ( um curso ), conceito 3 ( quatro cursos ) e conceito 5 ( três cursos ).
Na região Centro-Oeste, seis dos 12 cursos participantes eram de Instituições Privadas ( 50,0 \% em termos regionais, tal como na região Norte ).
Destes, todos foram alocados igualmente aos conceitos modais 2, 3 e 4 ( dois cursos, \textbf{cada} ).
Nenhum curso ficou sem conceito ( SC ) ou recebeu os conceitos 1 ou 5 nessa categoria.
Dos seis cursos de Instituições Públicas, quatro foram avaliados no conceito 5, conceito modal, seguido de dois cursos que receberam o conceito 2.
Nesta região, nenhum curso de IES Pública ficou sem conceito ( SC ) nessa categoria.
O Gráfico 5.2 apresenta a distribuição do Conceito ENADE segundo a Categoria Administrativa da IES.
Os cursos em IES Públicas ( linha azul ) apresentam uma poligonal mais à direita do que os em IES Privadas ( linha verde ) e consequentemente uma distribuição de Conceitos ENADE com valores maiores. 5.3 CONCEITOS POR ORGANIZAÇÃO ACADÊMICA E POR GRANDE REGIÃO Na Tabela 5.3 encontra -se a distribuição dos conceitos atribuídos aos cursos participantes do ENADE/2013 na Área de Medicina, por Organização Acadêmica, segundo as Grandes Regiões brasileiras.
Dos 166 cursos participantes, 113 eram oferecidos em Universidades, 38 em Faculdades e os demais 15 em Centros Universitários.
Esta distribuição corresponde a, respectivamente, 68,1 \%, 22,9 \% e 9,0 \% dos cursos.
De acordo com os dados apresentados, de todos os 20 cursos avaliados com conceito 5, 17 eram vinculados a Universidades.
Nenhum curso em Universidade ficou sem conceito ( SC ).
Esse tipo de Organização Acadêmica teve o conceito 3 como modal, com 40 cursos ( 35,4 \% ).
Os demais cursos de Universidades avaliados receberam os conceitos 1 ( seis cursos ), 2 ( 15 cursos ), 4 ( 35 cursos ) e conceito 5 ( 17 cursos, como já mencionado ).
Entre os cursos em Centros Universitários, os conceitos modais foram 3 e 4, ambos com cinco cursos, \textbf{cada}.
Neste tipo de Organização Acadêmica nenhum curso ficou sem conceito ( SC ) ou recebeu o conceito 5.
Os outros cursos neste tipo de Organização Acadêmica receberam os conceitos 1 ( três cursos ) e 2 ( dois cursos ).
Nas Faculdades, nenhum dos 38 cursos ficou sem conceito ( SC ), enquanto 14 ( 36,8 \% ) receberam o mesmo conceito modal dos demais tipos de Organização Acadêmica, conceito 3.
Dos demais cursos neste tipo de Organização Acadêmica, dois receberam conceito 1, 11 conceito 2, oito conceito 4 e três conceito 5.
Considerando -se separadamente as regiões brasileiras, verifica -se que, na região Norte, as Universidades concentraram sete dos 14 cursos participantes.
Dos cursos em Universidades, nenhum ficou sem conceito ( SC ) ou recebeu conceitos 2 ou 5.
Dois cursos ficaram alocados igualmente aos conceitos 1 e 4 ( um curso, \textbf{cada} ).
Cinco receberam o conceito modal, 3.
Os Centros Universitários da região Norte foram representados por dois cursos, destes, todos receberam somente o conceito 2.
Nesta categoria e região, nenhum curso ficou sem conceito ( SC ).
As Faculdades participaram com cinco cursos na região Norte.
Destes, dois cursos receberam conceito 1 ( 40,0 \% ), enquanto três cursos receberam conceito 2 ( 60,0 \%, conceito modal ).
Nenhum ficou sem conceito ( SC ) nessa categoria e região.
Na região Nordeste, as Universidades participaram com 24 dos 34 cursos na Área de Medicina da região.
Nenhum dos cursos em Universidades no Nordeste ficou sem conceito ( SC ).
O conceito modal foi 4, com 13 cursos.
Os demais receberam conceito 1 ( um curso ), 2 ( dois cursos ), 3 ( seis cursos ) e conceito 5 ( também, dois cursos ).
Os Centros Universitários contaram com dois cursos participantes na região Nordeste, os quais ficaram alocados igualmente aos conceitos 3 e 4 ( um curso, \textbf{cada} ), enquanto nenhum curso ficou sem conceito ( SC ) ou recebeu os demais conceitos restantes nessa categoria e região.
As Faculdades foram representadas por oito cursos na região Nordeste, dos quais nenhum ficou sem conceito ( SC ) ou recebeu os conceitos 1 ou 5.
O conceito modal foi 2, com cinco cursos.
Dos cursos restantes, dois receberam conceito 3 e somente um curso recebeu conceito 4.
Na região Sudeste, as Universidades concentraram 46 dos 77 cursos da região.
Entre os cursos em Universidades na região, o conceito modal foi 3 com 17 cursos, enquanto nenhum curso ficou sem conceito ( SC ).
Os demais cursos receberam os conceitos 1 ( três cursos ), 2 ( 12 cursos ) e 4 e 5(sete cursos, \textbf{cada} ).
Os Centros Universitários participaram com dez cursos na região Sudeste, dos quais quatro obtiveram conceito modal, 3, enquanto nenhum ficou sem conceito ( SC ) ou recebeu os conceitos 2 ou 5.
Seis cursos ficaram alocados igualmente aos conceitos 1 e 4 ( três cursos, \textbf{cada} ).
As Faculdades foram representadas por 21 cursos na região Sudeste, que se distribuíram nos conceitos 2 e 5 ( dois cursos, \textbf{cada} ), 3 ( dez cursos, conceito modal ) e 4 ( sete cursos ).
Nenhum curso ficou sem conceito ( SC ) ou recebeu o conceito 1.
Dos 29 cursos da região Sul, 27 eram de Universidades, para os quais o conceito modal foi 4, com 11 cursos.
Nesse tipo de organização, nenhum dos cursos ficou sem conceito ( SC ) ou recebeu conceito 2.
Os demais receberam os conceitos 1 ( um curso ), 3 ( dez cursos ) e 5 ( cinco cursos ).
Não houve nenhum curso vinculado a Organização Acadêmica Centro Universitário na região Sul.
Foram dois os cursos vinculados a Faculdades na região Sul, os quais receberam o conceito 3.
Nenhum curso ficou sem conceito ( SC ) ou recebeu os demais conceitos restantes nessa categoria e região.
Na região Centro-Oeste, nove dos 12 cursos eram de Universidades.
Nesse tipo de organização, nenhum curso ficou sem conceito ( SC ) ou recebeu o conceito 1.
Além disso, os conceitos modais foram 4 e 5, distribuídos igualmente em três cursos, \textbf{cada}.
Os outros cursos obtiveram os conceitos 2 ( um curso ) e 3 ( dois cursos ).
A Organização Acadêmica Centro Universitário da região Centro-Oeste contou somente com um único curso que recebeu o conceito 4.
Nenhum curso ficou sem conceito ( SC ) ou recebeu os demais conceitos restantes nessa categoria.
Dos dois cursos em Faculdades na região Centro-Oeste, ambos ficaram alocados igualmente aos conceitos 2 e 5 ( um curso, \textbf{cada} ).
Nenhum ficou sem conceito ( SC ) ou recebeu os demais conceitos restantes nessa categoria e região.
O Gráfico 5.3 apresenta a distribuição do Conceito ENADE segundo a Organização Acadêmica da IES.
Os cursos em Universidades ( linha azul ) apresentam uma poligonal mais à direita do que as demais, denotando uma melhor distribuição de Conceitos ENADE.
No outro extremo temos os cursos em Faculdades ( linha vermelha ) e Centros Universitários ( linha verde ) com a poligonal mais à esquerda e consequentemente uma distribuição de Conceitos ENADE com menores valores.

% matched lemmas: cada
\end{textsample}
